\section{AHSME 1976}
        \begin{enumerate}
        	\item (\textit{AHSME 1976}) If one minus the reciprocal of $(1-x)$ equals the reciprocal of $(1-x)$, then $x$ equals


 $\textbf{(A) }-2\qquad \textbf{(B) }-1\qquad \textbf{(C) }1/2\qquad \textbf{(D) }2\qquad  \textbf{(E) }3$

 \textbf{Answer: }B



\textbf{Solution}

The reciprocal of $(1-x)$ is $\frac{1}{1-x}$, so our equation is 
 \[1-\frac{1}{1-x}=\frac{1}{1-x},\] which is equivalent to $\frac{1}{1-x}=\frac{1}{2}$. So, $1-x=2$ and $x=-1\Rightarrow \textbf{(B)}$.~MathJams

 


 


	\item (\textit{AHSME 1976}) For how many real numbers $x$ is $\sqrt{-(x+1)^2}$ a real number?


 $\textbf{(A) }\text{none}\qquad \textbf{(B) }\text{one}\qquad \textbf{(C) }\text{two}\qquad\\ \textbf{(D) }\text{a finite number greater than two}\qquad \textbf{(E) }\infty$

 \textbf{Answer: }B



\textbf{Solution}

$\sqrt{-(x+1)^2}$ is a real number, if and only if $-(x+1)^2$ is nonnegative. Since $(x+1)^2$ is always nonnegative, $-(x+1)^2$ is nonnegative only when $-(x+1)^2=0$, or when $x=-1 \Rightarrow \textbf{(B)}$.~MathJams

 


 


	\item (\textit{AHSME 1976}) The sum of the distances from one vertex of a square with sides of length $2$ to the midpoints of each of the sides of the square is


 $\textbf{(A) }2\sqrt{5}\qquad \textbf{(B) }2+\sqrt{3}\qquad \textbf{(C) }2+2\sqrt{3}\qquad \textbf{(D) }2+\sqrt{5}\qquad \textbf{(E) }2+2\sqrt{5}$

 \textbf{Answer: }E



\textbf{Solution}

The lengths to the side are $1, \sqrt{2^2+1^2}, \sqrt{2^2+1^2}, 1$, respectively. Therefore, the sum is $\boxed{\textbf{(E) } 2+2\sqrt{5}}$.~MathJams

 


 


	\item (\textit{AHSME 1976}) Let a geometric progression with n terms have first term one, common ratio $r$ and sum $s$, where $r$ and $s$ are not zero. 
The sum of the geometric progression formed by replacing each term of the original progression by its reciprocal is


 $\textbf{(A) }\frac{1}{s}\qquad \textbf{(B) }\frac{1}{r^ns}\qquad \textbf{(C) }\frac{s}{r^{n-1}}\qquad \textbf{(D) }\frac{r^n}{s}\qquad \textbf{(E) } \frac{r^{n-1}}{s}$

 \textbf{Answer: }C



\textbf{Solution}

The sum of a geometric series with $n$ terms, initial term $a$, and ratio $r$ is $\frac{a(1-r^n)}{1-r}$. So, $s=\frac{(1-r^n)}{1-r}$. Our initial sequence is $1, r, r^2, \dots, r^n$, and replacing each terms with its reciprocal gives us the sequence $1, \frac{1}{r}, \frac{1}{r^2}, \dots, \frac{1}{r^n}$. The sum is now $\frac{1-(\frac{1}{r^n})}{1-\frac{1}{r}}=\frac{\frac{(1-r^n)}{r^n}}{\frac{1-r}{r}}=\frac{s}{r^{n-1}}\Rightarrow \textbf{(C)}$.~MathJams

 


	\item (\textit{AHSME 1976}) How many integers greater than $10$ and less than $100$, written in base-$10$ notation, 
are increased by $9$ when their digits are reversed?


 $\textbf{(A)}\ 0 \qquad \textbf{(B)}\ 1 \qquad \textbf{(C)}\ 8 \qquad \textbf{(D)}\ 9 \qquad \textbf{(E)}\ 10$

 \textbf{Answer: }C



\textbf{Solution}

Let our two digit number be $\overline{ab}$, where $a$ is the tens digit, and $b$ is the ones digit. So, $\overline{ab}=10a+b$. When we reverse our digits, it becomes $10b+a$. So, $10a+b+9=10b+a\implies a-b=1$. So, our numbers are $12, 23, 34, 45, 56, 67, 78, 89\Rightarrow \textbf{(C)}$.~MathJams

 


	\item (\textit{AHSME 1976}) If $c$ is a real number and the negative of one of the solutions of $x^2-3x+c=0$ is a solution of $x^2+3x-c=0$, then the solutions of $x^2-3x+c=0$ are


 $\textbf{(A) }1,~2\qquad \textbf{(B) }-1,~-2\qquad \textbf{(C) }0,~3\qquad \textbf{(D) }0,~-3\qquad  \textbf{(E) }\frac{3}{2},~\frac{3}{2}$

 \textbf{Answer: }C



\textbf{Solution}

We let the roots of the first equation be $r,s$ and the roots of the second equation be $s, -t$. By Vieta's Formulas, $r+s=3$ and $s-t=-3$, $rs=c$ and $-st=c$. So, $r=t$. Thus, $t+s=3$, $s-t=3$, so $t=0$, and $s=3\Rightarrow \textbf{(C)}$.~MathJams

 


	\item (\textit{AHSME 1976}) If $x$ is a real number, then the quantity $(1-|x|)(1+x)$ is positive if and only if


 $\textbf{(A) }|x|<1\qquad \textbf{(B) }|x|>1\qquad \textbf{(C) }x<-1\text{ or }-1<x<1\qquad\\ \textbf{(D) }x<1\qquad  \textbf{(E) }x<-1$

 \textbf{Answer: }C



\textbf{Solution}

We divide our solution into three cases: that of $x < -1$, that of $-1 < x < 1$, and that of $x > 1$. (When $x = -1$ or $x = 1$, the expression is zero, therefore not positive.)

 If $x < -1$, then the first \href{https://artofproblemsolving.com/wiki/index.php?title=Factor}{factor} is negative, and the second factor is also negative.

 If $-1 < x < 1$, then the first factor is positive, and the second factor is also positive.

 If $x > 1$, the first factor is negative, but the second factor is positive.

 Combining this with the rules for signs and multiplication, we find that the expression is positive when $x < -1$ or when $-1 < x < 1$, so our answer is $\boxed{\textbf{(C)}}$ and we are done.

 


	\item (\textit{AHSME 1976}) A point in the plane, both of whose rectangular coordinates are integers with absolute values less than or equal to four, 
is chosen at random, with all such points having an equal probability of being chosen. What is the probability that the distance 
from the point to the origin is at most two units?


 $\textbf{(A) }\frac{13}{81}\qquad \textbf{(B) }\frac{15}{81}\qquad \textbf{(C) }\frac{13}{64}\qquad \textbf{(D) }\frac{\pi}{16}\qquad \textbf{(E) }\text{the square of a rational number}$

 \textbf{Answer: }A



\textbf{Solution 1}

The rectangular region will be an 8x8 grid which is centered at the origin. Since the distance from the point chosen to the origin is at most two units, the favorable region will NOT be the $2 \times 2$ grid centered at the origin. It will be a circle of radius two units centered at the origin. The area of that circle is pi, and the area of the total region is 64. Therefore, the answer is \boxed{\textbf{(D) }\frac{\pi}{16}}

 


	\item (\textit{AHSME 1976}) In triangle $ABC$, $D$ is the midpoint of $AB$; $E$ is the midpoint of $DB$; and $F$ is the midpoint of $BC$. 
If the area of $\triangle ABC$ is $96$, then the area of $\triangle AEF$ is


 $\textbf{(A) }16\qquad \textbf{(B) }24\qquad \textbf{(C) }32\qquad \textbf{(D) }36\qquad  \textbf{(E) }48$

 \textbf{Answer: }D



\textbf{Solution}

Let $[ABC]$ denote the area of $ABC$Since $F$ is the midpoint of $|BC|$, we have $[ABF] = [ACF] = \frac{96}{2} = 48$$D$ is the midpoint of $|AB|$, so $[ADF] = [BDF] = \frac{48}{2} = 24$Now we just need to find $[DEF]$, since $[AEF] = [ADF] + [DEF]$In $\triangle DBF$, $E$ is the midpoint of $DB$, so $[DEF] = [EBF] = \frac{24}{2} = 12$Thus, 
 \[[AEF] = [ADF] + [DEF] = 24+12 = 36\]$\boxed{D}$



	\item (\textit{AHSME 1976}) If $m,~n,~p$, and $q$ are real numbers and $f(x)=mx+n$ and $g(x)=px+q$, then the equation $f(g(x))=g(f(x))$ has a solution


 $\textbf{(A) }\text{for all choices of }m,~n,~p, \text{ and } q\qquad\\ \textbf{(B) }\text{if and only if }m=p\text{ and }n=q\qquad\\ \textbf{(C) }\text{if and only if }mq-np=0\qquad\\ \textbf{(D) }\text{if and only if }n(1-p)-q(1-m)=0\qquad\\ \textbf{(E) }\text{if and only if }(1-n)(1-p)-(1-q)(1-m)=0$

 \textbf{Answer: }D



\textbf{Solution}


 \[f(g(x) = g(f(x)) \qquad \implies \qquad m(px + q) + n = p(mx + n) + q\]

 This simplifies to $mq + n = pn + q$ or $n(1 - p) - q(1 - m) = 0.$

 
The answer is $\boxed{\textbf{(D)}}.$



	\item (\textit{AHSME 1976}) Which of the following statements is (are) equivalent to the statement 
"If the pink elephant on planet alpha has purple eyes, then the wild pig on planet beta does not have a long nose"?


 $\begin{array}{l}\textbf{I. }\text{"If the wild pig on planet beta has a long nose, then the pink elephant on planet alpha has purple eyes."}\\ \textbf{II. }\text{"If the pink elephant on planet alpha does not have purple eyes, then the wild pig on planet beta does not have a long nose."}\\ \textbf{III. }\text{"If the wild pig on planet beta has a long nose, then the pink elephant on planet alpha does not have purple eyes."}\\ \textbf{IV. }\text{"The pink elephant on planet alpha does not have purple eyes, or the wild pig on planet beta does not have a long nose."}\end{array}$
$\textbf{(A) }\textbf{I. }\text{and }\textbf{II. }\text{only}\qquad \textbf{(B) }\textbf{III. }\text{and }\textbf{IV. }\text{only}\qquad \textbf{(C) }\textbf{II. }\text{and }\textbf{IV. }\text{only}\qquad \textbf{(D) }\textbf{II. }\text{and }\textbf{III. }\text{only}\qquad \textbf{(E) }\textbf{III. }\text{only}$

 \textbf{Answer: }B



\textbf{Solution}

The following statements are all logically equivalent.

 I. P implies Q (pink elephant on planet alpha has purple eyes ---> wild pig on planet beta does not have a long nose)

 
III. Not Q implies not P (contrapositive)

 IV. Not P nor Q (inverse)

 
I is the given statement, so III and IV work. The answer is $\boxed{\textbf{(B)}}.$



	\item (\textit{AHSME 1976}) A supermarket has $128$ crates of apples. Each crate contains at least $120$ apples and at most $144$ apples. 
What is the largest integer $n$ such that there must be at least $n$ crates containing the same number of apples?


 $\textbf{(A) }4\qquad \textbf{(B) }5\qquad \textbf{(C) }6\qquad \textbf{(D) }24\qquad  \textbf{(E) }25$

 \textbf{Answer: }C



\textbf{Solution}

To find the largest number of "repeated" crates necessary, we must account for all the possibilities of the number of apples in each crate. Since each crate contains a minimum of $120$ apples and a maximum of $144$ apples, there are $144 - 120 + 1 = 25$ different amounts possible for the number of apples per crate.

 Now, we have to count for the worst case scenario: the $25$ amounts are repeated as many times as possible.

 $25$ can go into $128$ exactly $5$ times because $5 \cdot 25 = 125$, which is less than $128$. This leaves a remainder of $3$ crates.

 The worst case scenario would be that these $3$ crates have a different number of apples each. It doesn't actually matter how many apples are in these $3$ crates because any of the $25$ values would be repeated again anyway. So, the answer is $5 + 1 = \boxed{\textbf{(C) }6}$ ~\href{https://artofproblemsolving.com/wiki/index.php/User:Jiang147369}{jiang147369}

 


	\item (\textit{AHSME 1976}) If $x$ cows give $x+1$ cans of milk in $x+2$ days, how many days will it take $x+3$ cows to give $x+5$ cans of milk?


 $\textbf{(A) }\frac{x(x+2)(x+5)}{(x+1)(x+3)}\qquad \textbf{(B) }\frac{x(x+1)(x+5)}{(x+2)(x+3)}\qquad\\ \textbf{(C) }\frac{(x+1)(x+3)(x+5)}{x(x+2)}\qquad \textbf{(D) }\frac{(x+1)(x+3)}{x(x+2)(x+5)}\qquad \\ \textbf{(E) }\text{none of these}$

 \textbf{Answer: }A



\textbf{Solution}

First, the problem states:

  $x$ cows $\qquad x+1$ cans $\qquad x+2$ days 

 Multiply the current number of cows by $\frac{x+3}{x}$ to get the number of cows you want, and divide the number of days by the same amount because an increase in cows will cause a decrease in time.

  $x \cdot \frac{x+3}{x} \qquad x+1 \qquad (x+2) \cdot \frac{x}{x+3}$

 Finally, multiply the current amount of cans by $\frac{x+5}{x+1}$ to get the number of cans you want, and multiply the number of days by the same amount because an increase in cans will cause an increase in time.

  $x \cdot \frac{x+3}{x} \qquad (x+1) \cdot \frac{x+5}{x+1} \qquad (x+2) \cdot \frac{x}{x+3} \cdot \frac{x+5}{x+1}$ 

 Therefore, our answer is $\boxed{\textbf{(A) }\frac{x(x+2)(x+5)}{(x+1)(x+3)}}$ ~\href{https://artofproblemsolving.com/wiki/index.php/User:Jiang147369}{jiang147369}

 


	\item (\textit{AHSME 1976}) The measures of the interior angles of a convex polygon are in arithmetic progression. 
If the smallest angle is $100^\circ$, and the largest is $140^\circ$, then the number of sides the polygon has is 


 $\textbf{(A) }6\qquad \textbf{(B) }8\qquad \textbf{(C) }10\qquad \textbf{(D) }11\qquad  \textbf{(E) }12$

 \textbf{Answer: }A



\textbf{Solution}

Let $n$ equal the number of sides the polygon has. The sum of all the interior angles of a polygon is: $180(n-2)$.

 The formula for an arithmetic series is $\frac{n(a_1 + a_n)}{2}$. Set this equal to $180(n-2)$ and solve. In this case, $a_1=100$ and $a_n=140$.

 Our equation becomes $\frac{n(100+140)}{2} = 180(n-2) \Rightarrow 240n = 360(n-2) \Rightarrow 120n = 720$.

 Simplifying, we get $n = \boxed{\textbf{(A) }6}$ ~\href{https://artofproblemsolving.com/wiki/index.php/User:Jiang147369}{jiang147369}

 


	\item (\textit{AHSME 1976}) If $r$ is the remainder when each of the numbers $1059,~1417$, and $2312$ is divided by $d$, where $d$ is an integer greater than $1$, then $d-r$ equals 


 $\textbf{(A) }1\qquad \textbf{(B) }15\qquad \textbf{(C) }179\qquad \textbf{(D) }d-15\qquad  \textbf{(E) }d-1$

 \textbf{Answer: }B



\textbf{Solution}

We are given these congruences:

 $1059 \equiv r \pmod{d} \qquad$ (i)

 $1417 \equiv r \pmod{d} \qquad$ (ii)

 $2312 \equiv r \pmod{d} \qquad$ (iii)

 
Let's make a new congruence by subtracting (i) from (ii), which results in 
 \[358 \equiv 0 \pmod{d}.\]Subtract (ii) from (iii) to get 
 \[895 \equiv 0 \pmod{d}.\]

 
Now we know that $358$ and $895$ are both multiples of $d$. Their prime factorizations are $358=2 \cdot 179$ and $895=5 \cdot 179$, so their common factor is $179$, which means $d=179$.

 
Plug $d=179$ back into any of the original congruences to get $r=164$. Then, $d-r=179-164= \boxed{\textbf{(B) }15}$. ~\href{https://artofproblemsolving.com/wiki/index.php?title=User:Jiang147369}{jiang147369}

 


 


	\item (\textit{AHSME 1976}) In triangles $ABC$ and $DEF$, lengths $AC,~BC,~DF$, and $EF$ are all equal. Length $AB$ is twice the length of the altitude of $\triangle DEF$
from $F$ to $DE$. Which of the following statements is (are) true?


 $\begin{array}{l}\textbf{I. }\angle ACB \text{ and }\angle DFE\text{ must be complementary.}\\ \textbf{II. }\angle ACB \text{ and }\angle DFE\text{ must be supplementary.}\\ \textbf{III. }\text{The area of }\triangle ABC\text{ must equal the area of }\triangle DEF.\\ \textbf{IV. }\text{The area of }\triangle ABC\text{ must equal twice the area of }\triangle DEF.\end{array}$
$\textbf{(A) }\textbf{II. }\text{only}\qquad \textbf{(B) }\textbf{III. }\text{only}\qquad \textbf{(C) }\textbf{IV. }\text{only}\qquad \textbf{(D) }\text{I. }\text{and }\textbf{III. }\text{only}\qquad  \textbf{(E) }\textbf{II. }\text{and }\textbf{III. }\text{only}$

 \textbf{Answer: }E



\textbf{Solution}

Let $AC=BC=DF=EF=x$ and let the altitude from $F$ to $\overline{DE}$ have length $y$. This gives that $AB=2y$. Since $\triangle ABC$ and $\triangle DEF$ are isosceles, we have that $\angle ACB = 2\sin ^{-1} \left(\frac{y}{x} \right)$ and $\angle DFE = 2\cos ^{-1} \left(\frac{y}{x} \right)$. Since the sum of the inverse sine and cosine of any possible sine/cosine fraction is $90^\circ$, we have that $m \angle ACB + m \angle DFE = 180^\circ$. Thus, $\angle ACB$ and $\angle DFE$ are supplementary, so $\mathbf{II.}$ is true while $\mathbf{I.}$ is not. 

 
Applying the Pythagorean Theorem gives that the length of the altitude from $C$ to $\overline{AB}$ is $\sqrt{x^2-y^2}$ and $DE = 2 \sqrt{x^2 - y^2}$. This means that $[ABC]=[DEF]=y \sqrt{x^2-y^2}.$ So we also have that $\mathbf{III.}$ is true while $\mathbf{IV.}$ is not. 

 
Thus, our answer is $\textbf{(E) }\textbf{II. }\text{and }\textbf{III. }\text{only.}$



	\item (\textit{AHSME 1976}) If $\theta$ is an acute angle, and $\sin 2\theta=a$, then $\sin\theta+\cos\theta$ equals


 $\textbf{(A) }\sqrt{a+1}\qquad \textbf{(B) }(\sqrt{2}-1)a+1\qquad \textbf{(C) }\sqrt{a+1}-\sqrt{a^2-a}\qquad\\ \textbf{(D) }\sqrt{a+1}+\sqrt{a^2-a}\qquad  \textbf{(E) }\sqrt{a+1}+a^2-a$

 \textbf{Answer: }A



\textbf{Solution}

Let $x = \sin\theta+\cos\theta$, so we want to find $x$. First, square the expression to get $x^2 = \sin^2 \theta + 2\sin\theta\cos\theta + \cos^2 \theta$. Recall that $\sin^2 \theta + \cos^2 \theta = 1$ and $\sin 2\theta = 2\sin\theta\cos\theta$. Plugging these in, the equation simplifies to $x^2 = 1 + \sin 2\theta$. Given that $\sin 2\theta=a$, the equation becomes $x^2=1+a$. Take the square root of both sides to get $x = \sqrt{1+a}$.

 Hence, the answer is $\boxed{\textbf{(A) }\sqrt{a+1}}$ ~\href{https://artofproblemsolving.com/wiki/index.php/User:Jiang147369}{jiang147369}

 


	\item (\textit{AHSME 1976}) 
                    \begin{figure}[H]
                        \centering
                        \includegraphics[width=0.5\textwidth]{Problem/AHSME_I_1976_18_Problem_Figure_1.png}
                    \end{figure}
                    


 In the adjoining figure, $AB$ is tangent at $A$ to the circle with center $O$; point $D$ is interior to the circle; 
and $DB$ intersects the circle at $C$. If $BC=DC=3$, $OD=2$, and $AB=6$, then the radius of the circle is


 $\textbf{(A) }3+\sqrt{3}\qquad \textbf{(B) }15/\pi\qquad \textbf{(C) }9/2\qquad \textbf{(D) }2\sqrt{6}\qquad  \textbf{(E) }\sqrt{22}$

 \textbf{Answer: }E



\textbf{Solution}

Extend $\overline{BD}$ until it touches the opposite side of the circle, say at point $E.$ By power of a point, we have $AB^2=(BC)(BE),$ so $BE=\frac{6^2}{3}=12.$ Therefore, $DE=12-3-3=6.$ 

 Now extend $\overline{OD}$ in both directions so that it intersects the circle in two points (in other words, draw the diameter of the circle that contains $OD$). Let the two endpoints of this diameter be $P$ and $Q,$ where $Q$ is closer to $C.$ Again use power of a point. We have $(PD)(DQ)=(CD)(DE)=3 \cdot 6=18.$ But if the radius of the circle is $r,$ we see that $PD=r+2$ and $DQ=r-2,$ so we have the equation $(r+2)(r-2)=18.$ Solving gives $r=\sqrt{22}.$



	\item (\textit{AHSME 1976}) A polynomial $p(x)$ has remainder three when divided by $x-1$ and remainder five when divided by $x-3$. 
The remainder when $p(x)$ is divided by $(x-1)(x-3)$ is


 $\textbf{(A) }x-2\qquad \textbf{(B) }x+2\qquad \textbf{(C) }2\qquad \textbf{(D) }8\qquad  \textbf{(E) }15$

 \textbf{Answer: }B



\textbf{Solution}

We know that $p(1)=3$ and $p(3)=5$. We write $p(x)$ as $p(x)=q(x)(x-1)(x-3)+r(x)$, where $r(x)=ax+b$. Plugging in $x=1$, we get $a+b=3$. Plugging in $x=3$, we know $3a+b=5$. We have a systems of equations, where we can solve that $a=1$ and $b=2$. So, our answer is $1(x)+(2)=x+2\Rightarrow \textbf{(B)}$. ~MathJams

 \textbf{1976 AHSME Problems} 


	\item (\textit{AHSME 1976}) Let $a,~b$, and $x$ be positive real numbers distinct from one. Then $4(\log_ax)^2+3(\log_bx)^2=8(\log_ax)(\log_bx)$
$\textbf{(A) }\text{for all values of }a,~b,\text{ and }x\qquad\\ \textbf{(B) }\text{if and only if }a=b^2\qquad\\ \textbf{(C) }\text{if and only if }b=a^2\qquad\\ \textbf{(D) }\text{if and only if }x=ab\qquad\\ \textbf{(E) }\text{for none of these}$

 \textbf{Answer: }E



\textbf{Solution}

Because $\log_{m} n = \dfrac{\log n}{\log m}$, $4(\log_{a} x)^2+3(\log_{b} x)^2 =$ $\dfrac{4(\log x)^2}{(\log a)^2}+\dfrac{3(\log x)^2}{(\log b)^2} = \dfrac{(\log x)^2(4(\log a)^2+3(\log b)^2)}{(\log a \log b)^2}$.

 Therefore,

 
 \[\dfrac{(\log x)^2(4(\log a)^2+3(\log b)^2)}{(\log a \log b)^2} = \frac{8(\log x)^2}{\log a \log b}\]
 \[\dfrac{4(\log a)^2+3(\log b)^2}{\log a \log b} = 8\]

 \[\dfrac{4(\log a)^2}{\log a \log b} + \dfrac{3(\log b)^2}{\log a \log b} = 8\]

 \[4 \log_{b} a + 3 \log_{a} b = 8\]

 Now let $k = \log_{b} a$. Our equation becomes

 
 \[4k + \frac{3}{k} = 8\]
 \[4k^2 - 8k + 3 = 0\]
 \[k = \frac{8 \pm \sqrt{64-4(4)(3)}}{8}\]
 \[k = \frac{1}{2}, \frac{3}{2}\]

 Therefore, either $b = a^2$ or $b = \sqrt[3] {a^2}$. Since no option matches both of these solutions, the answer is $\boxed{\textbf{(E)}}$.

 - mako17

 


	\item (\textit{AHSME 1976}) What is the smallest positive odd integer $n$ such that the product $2^{1/7}2^{3/7}\cdots2^{(2n+1)/7}$ is greater than $1000$? 
(In the product the denominators of the exponents are all sevens, and the numerators are the successive odd integers from $1$ to $2n+1$.)


 $\textbf{(A) }7\qquad \textbf{(B) }9\qquad \textbf{(C) }11\qquad \textbf{(D) }17\qquad  \textbf{(E) }19$

 \textbf{Answer: }B



\textbf{Solution}

Combine the terms in the product to get $2^{\frac{1+3+5+ \dots +(2n-1)+(2n+1)}{7}}$.

 The exponent can be simplified to 
 \[\frac{1+3+5+ \dots +(2n-1)+(2n+1)}{7} \Rightarrow \frac{\frac{n(1+(2n+1))}{2}}{7} \Rightarrow \frac{n^2}{7}.\]

 We want this inequality to be true with the smallest positive odd integer value of $n$: 
 \[2^{\frac{n^2}{7}} > 1000.\]

 Now, let's test the answer choices. For $n=7$, we have $2^{49/7}=2^{7}<1000$. For $n=9$, we have $2^{81/7}>2^{10}>1000$.

 
So our answer is $\boxed{\textbf{(B) }9}$. ~\href{https://artofproblemsolving.com/wiki/index.php?title=User:Jiang147369}{jiang147369}

 


 


	\item (\textit{AHSME 1976}) Given an equilateral triangle with side of length $s$, consider the locus of all points $\mathit{P}$ in the plane of the 
triangle such that the sum of the squares of the distances from $\mathit{P}$ to the vertices of the triangle is a fixed number $a$. This locus


 $\textbf{(A) }\text{is a circle if }a>s^2\qquad\\ \textbf{(B) }\text{contains only three points if }a=2s^2\text{ and is a circle if }a>2s^2\qquad\\ \textbf{(C) }\text{is a circle with positive radius only if }s^2<a<2s^2\qquad\\ \textbf{(D) }\text{contains only a finite number of points for any value of }a\qquad\\ \textbf{(E) }\text{is none of these}$

 \textbf{Answer: }A



\textbf{Solution}

If we define point $\mathit{P} = (x, y)$, and the three points of the triangle as $P_1 = (x_1,y_1)$, $P_2 = (x_2,y_2)$, and $P_3 = (x_3,y_3)$, then we have two equations to work with.  

 First, $\mathit{P}$ must satisfy the equation:

 $|P - P_1|^2 + |P - P_2|^2 + |P - P_3|^2 =\newline [(x-x_1)^2+(y-y_1)^2]+[(x-x_2)^2+(y-y_2)^2]+[(x-x_3)^2+(y-y_3)^2] = a$

 Second, the fact that the three sides of the triangle all have length $s$ gives us:$|P_2 - P_1|^2 + |P_3 - P_2|^2 + |P_1 - P_3|^2 =\newline [(x_2-x_1)^2+(y_2-y_1)^2]+[(x_3-x_2)^2+(y_3-y_2)^2]+[(x_1-x_3)^2+(y_1-y_3)^2] = 3s^2$

 Define the center of the triangle as $P_C = (x_C, y_C)$.  

 Note that $x_C = \frac{x_1+x_2+x_3}{3}$ and $y_C = \frac{y_1+y_2+y_3}{3}$.  

 Also note that the distance between $P_C$ and each of the triangle vertices is $\frac{\sqrt{3}s}{3}$, so:

 $|P_C - P_1|^2 + |P_C - P_2|^2 + |P_C - P_3|^2 =\newline [(x_C-x_1)^2+(y_C-y_1)^2]+[(x_C-x_2)^2+(y_C-y_2)^2]+[(x_C-x_3)^2+(y_C-y_3)^2] = s^2$

 Rewrite $(x-x_1) = (x - x_C) + (x_C - x_1)$, and similarly for $(y-y_1), (x-x_2)$, etc.

 Now expanding the first sum and using the second sum, we get:

 $a = 3[(x-x_C)^2+(y-y_C)^2] + s^2\newline  + 2(x - x_C)[(x_C - x_1) + (x_C - x_2) + (x_C - x_3)] + 2(y - y_C)[(y_C - y_1) + (y_C - y_2) + (y_C - y_3)]$

 Because $P_C$ is centroid of the triangle, its coordinates are the averages of the coordinates of the three points.  Thus, $3x_C = x_1 + x_2 + x_3$, and the latter parts of the above equations are zero.

 Thus, $a - s^2 = 3[(x-x_C)^2+(y-y_C)^2]$

 This means that when $a > s^2$, the locus of points $P$ will be a circle centered at $P_C$.

 Therefore, the correct answer is $\boxed{\textbf{(A)}}$.

 This problem could be also solved with coordinate geometry, where we place $P_1$ at the origin and $P_2$ along the positive x-axis, or even let $P_C$ be the origin and place $P_1$ along the negative x-axis.

 


	\item (\textit{AHSME 1976}) For integers $k$ and $n$ such that $1\le k<n$, let $C^n_k=\frac{n!}{k!(n-k)!}$. Then $\left(\frac{n-2k-1}{k+1}\right)C^n_k$ is an integer


 $\textbf{(A) }\text{for all }k\text{ and }n\qquad \\ \textbf{(B) }\text{for all even values of }k\text{ and }n,\text{ but not for all }k\text{ and }n\qquad \\ \textbf{(C) }\text{for all odd values of }k\text{ and }n,\text{ but not for all }k\text{ and }n\qquad \\ \textbf{(D) }\text{if }k=1\text{ or }n-1,\text{ but not for all odd values }k\text{ and }n\qquad \\ \textbf{(E) }\text{if }n\text{ is divisible by }k,\text{ but not for all even values }k\text{ and }n$

 \textbf{Answer: }A



\textbf{Solution}

We know $C^n_k = \binom{n}{k}$, so let's rewrite the expression as $\left(\frac{n-2k-1}{k+1}\right) \binom{n}{k}$. Notice that 
 \[n-2k-1 = n-2(k+1)+1 = (n+1)-2(k+1).\]

 This allows us to rewrite the expression as 
 \[\left(\frac{(n+1)-2(k+1)}{k+1}\right) \binom{n}{k}.\]

 From here, we just have to do some algebra to get
 \[\left(\frac{(n+1)-2(k+1)}{k+1}\right) \binom{n}{k} = \left( \frac{n+1}{k+1}-2 \right) \frac{n!}{k!(n-k)!}\]
 \[= \frac{(n+1)!}{(k+1)!(n-k)!} - 2 \cdot \frac{n!}{k!(n-k)!}\]
 \[= \binom{n+1}{k+1}-2\binom{n}{k},\]so $\left(\frac{n-2k-1}{k+1}\right)C^n_k$ is an integer $\boxed{\textbf{(A) }\text{for all }k\text{ and }n}$. ~\href{https://artofproblemsolving.com/wiki/index.php?title=User:Jiang147369}{jiang147369}

 


 


	\item (\textit{AHSME 1976}) In the adjoining figure, circle $K$ has diameter $AB$; circle $L$ is tangent to circle $K$ and to $AB$ at the center of circle $K$; and circle $M$ tangent to circle $K$, to circle $L$ and $AB$. The ratio of the area of circle $K$ to the area of circle $M$ is

                    \begin{figure}[H]
                        \centering
                        \includegraphics[width=0.5\textwidth]{Problem/AHSME_I_1976_24_Problem_Figure_1.png}
                    \end{figure}
                    
$\textbf{(A) }12\qquad \textbf{(B) }14\qquad \textbf{(C) }16\qquad \textbf{(D) }18\qquad  \textbf{(E) }\text{not an integer}$

 \textbf{Answer: }C



\textbf{Solution}

Let $R$ and $r$ be the radius of $\odot K$ and the radius of $\odot M,$ respectively. It follows that the radius of $\odot L$ is $\frac{R}{2}.$

 Suppose $P$ is the foot of the perpendicular from $M$ to $\overline{KL}.$ We construct the auxiliary lines, as shown below:
                    \begin{figure}[H]
                        \centering
                        \includegraphics[width=0.5\textwidth]{Solution/AHSME_I_1976_24_Solution_Figure_1.png}
                    \end{figure}
                    In right $\triangle KPM,$ we have $KP=r$ and $KM=R-r.$ By the Pythagorean Theorem, we get $PM^2=(R-r)^2-r^2.$

 In right $\triangle LPM,$ we have $LP=\frac{R}{2}-r$ and $LM=\frac{R}{2}+r.$ By the Pythagorean Theorem, we get $PM^2=\left(\frac{R}{2}+r\right)^2-\left(\frac{R}{2}-r\right)^2.$

 We equate the expressions for $PM^2,$ then simplify:
 \begin{align*} (R-r)^2-r^2&=\left(\frac{R}{2}+r\right)^2-\left(\frac{R}{2}-r\right)^2 \\ \left(R^2-2Rr+r^2\right)-r^2&=\left(\frac{R^2}{4}+Rr+r^2\right)-\left(\frac{R^2}{4}-Rr+r^2\right) \\ R^2-2Rr&=2Rr \\ R^2&=4Rr \\ R&=4r. \end{align*}
Therefore, the ratio of the area of $\odot K$ to the area of $\odot M$ is $\frac{\pi R^2}{\pi r^2}=\left(\frac{R}{r}\right)^2=\boxed{\textbf{(C) }16}.$



	\item (\textit{AHSME 1976}) For a sequence $u_1,u_2\dots$, define $\Delta^1(u_n)=u_{n+1}-u_n$ and, for all integer 
$k>1, \Delta^k(u_n)=\Delta^1(\Delta^{k-1}(u_n))$. 
If $u_n=n^3+n$, then $\Delta^k(u_n)=0$ for all $n$
$\textbf{(A) }\text{if }k=1\qquad \\ \textbf{(B) }\text{if }k=2,\text{ but not if }k=1\qquad \\ \textbf{(C) }\text{if }k=3,\text{ but not if }k=2\qquad \\ \textbf{(D) }\text{if }k=4,\text{ but not if }k=3\qquad\\ \textbf{(E) }\text{for no value of }k$

 \textbf{Answer: }D



\textbf{Solution}

Note that
 \begin{align*} \Delta^1(u_n) &= \left[(n+1)^3+(n+1)\right]-\left[n^3+n\right] \\ &= \left[n^3+3n^2+4n+2\right]-\left[n^3+n\right] \\ &= 3n^2+3n+2, \\ \Delta^2(u_n) &= \left[3(n+1)^2+3(n+1)+2\right]-\left[3n^2+3n+2\right] \\ &= \left[3n^2+9n+8\right]-\left[3n^2+3n+2\right] \\ &= 6n+6, \\ \Delta^3(u_n) &= \left[6(n+1)+6\right]-\left[6n+6\right] \\ &= \left[6n+12\right]-\left[6n+6\right] \\ &= 6, \\ \Delta^4(u_n) &= 6-6 \\ &= 0. \end{align*}
Therefore, the answer is $\boxed{\textbf{(D)}}.$

 More generally, we have $\Delta^k(u_n)=0$ for all $n$ if $k\geq4.$



	\item (\textit{AHSME 1976}) 
                    \begin{figure}[H]
                        \centering
                        \includegraphics[width=0.5\textwidth]{Problem/AHSME_I_1976_26_Problem_Figure_1.png}
                    \end{figure}
                    


 In the adjoining figure, every point of circle $\mathit{O'}$ is exterior to circle $\mathit{O}$. 
Let $\mathit{P}$ and $\mathit{Q}$ be the points of intersection of an internal common tangent with the two external common tangents. 
Then the length of $PQ$ is


 $\textbf{(A) }\text{the average of the lengths of the internal and external common tangents}\qquad\\ \textbf{(B) }\text{equal to the length of an external common tangent if and only if circles }\mathit{O}\text{ and }\mathit{O'} \text{have equal radii}\\ \textbf{(C) }\text{always equal to the length of an external common tangent}\qquad\\ \textbf{(D) }\text{greater than the length of an external common tangent}\qquad\\ \textbf{(E) }\text{the geometric mean of the lengths of the internal and external common tangents}$

 \textbf{Answer: }C



\textbf{Solution}

By the \href{https://artofproblemsolving.com/wiki/index.php?title=Two\_Tangent\_Theorem}{Two Tangent Theorem}, the segments drawn from point $P$ tangent to circle $O$ are congruent. Let them have length $p$. Similarly, the two segments drawn from point $Q$ tangent to circle $O^{\prime}$ are congruent. Let them have length $q$.

 Let the length of the two external common tangents be $x$. We know that the two segments drawn from point $P$ tangent to circle $O^{\prime}$ are congruent. One of these tangents (the one on the external common tangent) has length $x-p$, so the other (the one on the internal common tangent) also has length $x-p$ Doing the same thing for point $Q$ and circle $O$ shows that the segments drawn from $Q$ tangent to circle $O$ has length $x-q$. Thus, $PQ=\frac{x-p+x-q+p+q}2=\frac{2x}2=x$, because the numerator of this fraction counts the length of $\overline{PQ}$ twice. Because $x$ is the length of the external common tangent, our answer is $\boxed{\textbf{(C) }\text{always equal to the length of an external common tangent}}$.

 


	\item (\textit{AHSME 1976}) If 
 \[N=\frac{\sqrt{\sqrt{5}+2}+\sqrt{\sqrt{5}-2}}{\sqrt{\sqrt{5}+1}}-\sqrt{3-2\sqrt{2}},\] then $N$ equals


 $\textbf{(A) }1\qquad \textbf{(B) }2\sqrt{2}-1\qquad \textbf{(C) }\frac{\sqrt{5}}{2}\qquad \textbf{(D) }\sqrt{\frac{5}{2}}\qquad \textbf{(E) }\text{none of these}$

 \textbf{Answer: }A



\textbf{Solution 1}

Let $x=\frac{\sqrt{\sqrt{5}+2}+\sqrt{\sqrt{5}-2}}{\sqrt{\sqrt{5}+1}}$ and $y=\sqrt{3-2\sqrt{2}}.$

 Note that
 \begin{align*} x^2&=\frac{\left(\sqrt{\sqrt{5}+2}+\sqrt{\sqrt{5}-2}\right)^2}{\left(\sqrt{\sqrt{5}+1}\right)^2} \\ &=\frac{\left(\sqrt{5}+2\right)+2\cdot\left(\sqrt{\sqrt{5}+2}\cdot\sqrt{\sqrt{5}-2}\right)+\left(\sqrt{5}-2\right)}{\sqrt{5}+1} \\ &=\frac{\left(\sqrt{5}+2\right)+2+\left(\sqrt{5}-2\right)}{\sqrt{5}+1} \\ &=\frac{2\sqrt{5}+2}{\sqrt{5}+1} \\ &=2. \end{align*}
Since $x>0,$ we have $x=\sqrt{2}.$

 On the other hand, note that
 \begin{align*} y^2&=3-2\sqrt{2} \\ &=2-2\sqrt{2}+1 \\ &=\left(\sqrt{2}-1\right)^2. \end{align*}
Since $y>0,$ we have $y=\sqrt{2}-1.$

 Finally, the answer is 
 \[N=x-y=\boxed{\textbf{(A) }1}.\]




\textbf{Solution 2}

Let $x=\frac{\sqrt{\sqrt{5}+2}+\sqrt{\sqrt{5}-2}}{\sqrt{\sqrt{5}+1}}$ and $y=\sqrt{3-2\sqrt{2}}.$

 Note that
 \[x=\frac{\sqrt{\sqrt{5}+2}+\sqrt{\sqrt{5}-2}}{\sqrt{\sqrt{5}+1}}\cdot\frac{\sqrt{\sqrt{5}-1}}{\sqrt{\sqrt{5}-1}}=\frac{\sqrt{3+\sqrt{5}}+\sqrt{7-3\sqrt{5}}}{2}. \hspace{15mm} (\bigstar)\]We rewrite each term in the numerator separately:

 
            \begin{enumerate} 
        \item Let $\sqrt{a}+\sqrt{b}=\sqrt{3+\sqrt{5}}$ for some nonnegative rational numbers $a$ and $b.$ We square both sides of this equation, then rearrange:
 \begin{align*} a+b+2\sqrt{ab}&=3+\sqrt{5} \\ a+b+\sqrt{4ab}&=3+\sqrt{5}. \end{align*}
It follows that
 \begin{align*} a+b&=3, \\ ab&=\frac54. \end{align*}
By inspection, we get $\{a,b\}=\left\{\frac12,\frac52\right\}.$ Alternatively, we conclude that $a$ and $b$ are the solutions to the quadratic equation $t^2-3t+\frac54=0$ by Vieta's Formulas, in which $t=\frac12,\frac52.$  Therefore, we obtain 
 \[\sqrt{3+\sqrt{5}}=\sqrt{\frac12}+\sqrt{\frac52}=\frac{\sqrt2}{2}+\frac{\sqrt{10}}{2}.\]


        \item Similarly, we obtain 
 \[\sqrt{7-3\sqrt{5}}=\frac{3\sqrt2}{2}-\frac{\sqrt{10}}{2}.\]

            \end{enumerate}
            Substituting these results into $(\bigstar),$ we have 
 \[x=\frac{\left(\frac{\sqrt2}{2}+\frac{\sqrt{10}}{2}\right)+\left(\frac{3\sqrt2}{2}-\frac{\sqrt{10}}{2}\right)}{2}=\sqrt2.\]On the other hand, we have 
 \[y=\sqrt2-1\] by the argument of either Solution 1 or Solution 2.

 Finally, the answer is 
 \[N=x-y=\boxed{\textbf{(A) }1}.\]



	\item (\textit{AHSME 1976}) Lines $L_1,L_2,\dots,L_{100}$ are distinct. All lines $L_{4n}, n$ a positive integer, are parallel to each other. 
All lines $L_{4n-3}, n$ a positive integer, pass through a given point $A.$ The maximum number of points of intersection of pairs of lines from the complete set $\{L_1,L_2,\dots,L_{100}\}$ is


 $\textbf{(A) }4350\qquad \textbf{(B) }4351\qquad \textbf{(C) }4900\qquad \textbf{(D) }4901\qquad  \textbf{(E) }9851$

 \textbf{Answer: }B



\textbf{Solution}

We partition $\{L_1,L_2,\dots,L_{100}\}$ into three sets. Let
 \begin{align*} X &= \{L_n\mid n\equiv0\pmod{4}\}, \\ Y &= \{L_n\mid n\equiv1\pmod{4}\}, \\ Z &= \{L_n\mid n\equiv2,3\pmod{4}\}, \\ \end{align*}
from which $|X|=|Y|=25$ and $|Z|=50.$

 Any two distinct lines can intersect at most once. To maximize the number of points of intersection, note that each point must be passed by exactly two lines. If three or more lines pass through the same point, then we can create more points of intersection by translating the lines.

 We construct the sets one by one:

 
            \begin{enumerate} 
        \item We construct all lines in set $X.$ Since the lines in set $X$ are parallel to each other, they have $0$ points of intersection.

  
        \item We construct all lines in set $Y.$ The lines in set $Y$ have $1$ point of intersection with each other, namely $A.$ 

 Moreover, each line in set $Y$ can intersect each line in set $X$ once. So, there are $25^2=625$ points of intersection. 

 \textbf{Now, we have $\boldsymbol{1+625=626}$ additional points of intersection.} 

 

  
        \item We construct all lines in set $Z.$ The lines in set $Z$ can have $\binom{50}{2}=1225$ points of intersection with each other. 

 Moreover, each line in set $Z$ can intersect each line in sets $X$ and $Y$ once. So, there are $50^2=2500$ points of intersection. 

 \textbf{Now, we have $\boldsymbol{1225+2500=3725}$ additional points of intersection.} 

 

 
            \end{enumerate}
            Together, the answer is $626+3725=\boxed{\textbf{(B) }4351}.$



	\item (\textit{AHSME 1976}) Ann and Barbara were comparing their ages and found that Barbara is as old as Ann was when Barbara was as old as 
Ann had been when Barbara was half as old as Ann is. If the sum of their present ages is $44$ years, then Ann's age is


 $\textbf{(A) }22\qquad \textbf{(B) }24\qquad \textbf{(C) }25\qquad \textbf{(D) }26\qquad  \textbf{(E) }28$

 \textbf{Answer: }B



\textbf{Solution}

This problem is very wordy. Nonetheless, let $a$ and $b$ be Ann and Barbara's current ages, respectively. We are given that $a+b=44$. Let $y$ equal the difference between their ages, so $y=a-b$. Know that $y$ is constant because the difference between their ages will always be the same.

 Now, let's tackle the equation: $b=$ Ann's age when Barbara was Ann's age when Barbara was $\frac{a}{2}$. When Barbara was $\frac{a}{2}$ years old, Ann was $\frac{a}{2}+y$ years old. So the equation becomes $b=$ Ann's age when Barbara was $\frac{a}{2}+y$. Adding on their age difference again, we get $b = \frac{a}{2} + y + y \Rightarrow b = \frac{a}{2} + 2y$. Substitute $a-b$ back in for $y$ to get $b = \frac{a}{2} + 2(a-b)$. Simplify: $2b = a + 4(a-b) \Rightarrow 6b = 5a$. Solving $b$ in terms of $a$, we have $b = \frac{5a}{6}$. Substitute that back into the first equation of $a+b=44$ to get $\frac{11a}{6}=44$. Solve for $a$, and the answer is $\boxed{\textbf{(B) }24}$. ~\href{https://artofproblemsolving.com/wiki/index.php/User:Jiang147369}{jiang147369}

 


	\item (\textit{AHSME 1976}) How many distinct ordered triples $(x,y,z)$ satisfy the following equations?

 \begin{align*} x + 2y + 4z &= 12 \\ xy + 4yz + 2xz &= 22 \\ xyz &= 6 \end{align*}

$\textbf{(A) }\text{none}\qquad \textbf{(B) }1\qquad \textbf{(C) }2\qquad \textbf{(D) }4\qquad \textbf{(E) }6$

 \textbf{Answer: }E



\textbf{Solution}

The first equation suggests the substitution $(a,b,c)=(x,2y,4z),$ from which $(x,y,z)=\left(a,\frac b2,\frac c4\right).$

 We rewrite the given equations in terms of $a,b,$ and $c:$
 \begin{align*} a + b + c &= 12, \\ \frac{ab}{2} + \frac{bc}{2} + \frac{ac}{2} &= 22, \\ \frac{abc}{8} &= 6. \end{align*}
We clear fractions in these equations:
 \begin{align*} a + b + c &= 12, \\ ab + ac + bc &= 44, \\ abc &= 48. \end{align*}
By Vieta's formulas, note that $a,b,$ and $c$ are the roots of the equation 
 \[r^3 - 12r^2 + 44r - 48 = 0,\]which factors as
 \[(r - 2)(r - 4)(r - 6) = 0.\]It follows that $\{a,b,c\}=\{2,4,6\}.$ For each permutation of $(a,b,c),$ we find the corresponding ordered triple $(x,y,z)=\left(a,\frac b2,\frac c4\right),$ as shown below:
 \[\begin{array}{c|c|c||c|c|c} & & & & & \\ [-2.5ex] \boldsymbol{a} & \boldsymbol{b} & \boldsymbol{c} & \boldsymbol{x} & \boldsymbol{y} & \boldsymbol{z} \\ [0.5ex] \hline & & & & & \\ [-2ex] 2 & 4 & 6 & 2 & 2 & 3/2 \\ 2 & 6 & 4 & 2 & 3 & 1 \\ 4 & 2 & 6 & 4 & 1 & 3/2 \\ 4 & 6 & 2 & 4 & 3 & 1/2 \\ 6 & 2 & 4 & 6 & 1 & 1 \\ 6 & 4 & 2 & 6 & 2 & 1/2 \end{array}\]So, there are $\boxed{\textbf{(E) }6}$ such ordered triples $(x,y,z).$



\end{enumerate}